

% ══════════════════════════════════════════════════════════


% ── Numérotation à partir du chapitre 1 ───────────────────
\setcounter{chapter}{0}

% ══════════════════════════════════════════════════════════
\chapter{DevSecOps et Pratiques de Sécurité}
% ══════════════════════════════════════════════════════════

% ------------------------------------------------------------
%  CHAPEAU INTRODUCTIF
% ------------------------------------------------------------
L'introduction générale a mis en évidence deux constats fondamentaux~:
d'une part, les architectures microservices démultiplient la surface d'attaque
et rendent les communications interservices difficiles à sécuriser~; d'autre
part, le Service Mesh s'impose comme une infrastructure complémentaire capable
d'adresser les angles morts des approches de test classiques.

Ce chapitre construit le socle méthodologique qui structure l'ensemble du
projet. Il présente successivement le paradigme DevSecOps
(section~\ref{sec:devsecops}), les quatre grandes
familles de tests de sécurité applicative SAST, DAST, IAST et RASP
en mettant en évidence leurs complémentarités et leurs limites respectives
(section~\ref{sec:tests}), puis les enjeux spécifiques du DAST dans les
architectures microservices, dont les limites motivent directement l'approche
hybride proposée (section~\ref{sec:dast-micro}).

% ------------------------------------------------------------
%  1.1  OBJECTIFS DU PROJET
% ------------------------------------------------------------
\section{Objectifs du projet}
\label{sec:objectifs}

Sur la base de ces constats, ce projet vise à concevoir et valider une
approche de sécurité hybride intégrant le Service Mesh dans une démarche
DevSecOps, afin de couvrir l'ensemble des dimensions de sécurité d'une
architecture microservices aussi bien les flux externes que les
communications interservices internes.

Les objectifs poursuivis s'articulent autour de trois axes~:

\paragraph{Axe~1  Cadre méthodologique DevSecOps.}
Maîtriser le paradigme DevSecOps et les quatre
familles de tests de sécurité SAST, DAST, IAST, RASP afin
d'identifier les outils les mieux adaptés à chaque phase du cycle de
développement, et de caractériser précisément les limites du DAST
conventionnel face à la dynamicité des architectures microservices.

\paragraph{Axe~2  Architecture de sécurité proposée.}
Concevoir une architecture hybride combinant Istio Service Mesh et un pipeline
DAST automatisé (OWASP ZAP), couvrant à la fois les flux \emph{north-south}
(trafic externe) et \emph{east-west} (communications interservices internes),
et intégrée dans une pipeline CI/CD reproductible déclenchée à chaque
déploiement en environnement de staging.

\paragraph{Axe~3  Validation expérimentale.}
Déployer un prototype représentatif un système de paiement
microservices sur un cluster Kubernetes réel (Minikube avec Istio),
évaluer l'apport concret du Service Mesh face à douze scénarios d'attaque
couvrant les six catégories STRIDE, et mesurer les limites restantes de
l'approche.


  \textbf{Architecture proposée.}
  L'approche développée dans ce projet repose sur trois piliers
  complémentaires~: \textbf{(1)}~le Service Mesh (Istio) comme infrastructure
  de sécurité transparente assurant mTLS, contrôle d'accès et
  observabilité~; \textbf{(2)}~le DAST interne (OWASP ZAP) comme mécanisme
  d'évaluation continue de la posture de sécurité, étendu aux flux
  interservices~; \textbf{(3)}~la pipeline DevSecOps comme cadre
  d'intégration et d'automatisation de l'ensemble.


% ------------------------------------------------------------
%  1.2  INTRODUCTION AU DEVSECOPS
% ------------------------------------------------------------
\section{Introduction au DevSecOps}
\label{sec:devsecops}

\subsection{Du DevOps au DevSecOps : genèse et évolution}

Le mouvement DevOps, apparu au tournant des années 2010, visait à supprimer
le cloisonnement traditionnel entre les équipes de développement
(\emph{Development}) et d'exploitation (\emph{Operations}). En automatisant
les pipelines d'intégration et de déploiement continus (CI/CD), il a permis
d'accélérer considérablement la livraison logicielle tout en améliorant la
stabilité des systèmes.

Cependant, cette accélération a mis en évidence une faille structurelle~: la
sécurité restait souvent une préoccupation tardive, traitée en phase finale
par des équipes spécialisées travaillant en silo. Ce modèle, désigné par
l'expression \og{}Security as a Gatekeeper\fg{}, se révèle incompatible avec
des rythmes de déploiement quotidiens, car il crée des goulots d'étranglement
et laisse des vulnérabilités non détectées atteindre la production.

Le concept de DevSecOps- contraction de \emph{Development}, \emph{Security}
et \emph{Operations}~ émerge comme réponse à cette inadéquation. Il s'agit
d'une culture, d'un ensemble de pratiques et d'outils visant à intégrer la
sécurité à chaque étape du cycle de vie logiciel, en partageant la
responsabilité de la sécurité entre toutes les équipes.


  Le DevSecOps est une approche organisationnelle et technique qui intègre la
  sécurité comme une composante intrinsèque et continue du cycle de
  développement et d'exploitation logicielle, par opposition à un contrôle
  ponctuel en fin de cycle.


\subsection{Les principes fondateurs du DevSecOps}

Le DevSecOps repose sur quatre piliers interdépendants~:

\begin{enumerate}[leftmargin=*, label=\textbf{\arabic*.}]

  \item \textbf{Automatisation de la sécurité.} Les contrôles de sécurité
    (analyse statique du code, scan de dépendances, tests dynamiques) sont
    intégrés directement dans les pipelines CI/CD et s'exécutent à chaque
    commit, sans intervention manuelle systématique.

  \item \textbf{Collaboration et responsabilité partagée.} Les développeurs,
    les ingénieurs DevOps et les équipes de sécurité (SecOps) travaillent
    conjointement dès la phase de conception. La sécurité n'est plus la
    propriété exclusive d'une équipe spécialisée.

  \item \textbf{Observabilité et retour d'information rapide.} Des mécanismes
    de monitoring, de logging et d'alerting permettent de détecter et de
    corriger rapidement les incidents de sécurité en production.

  \item \textbf{Infrastructure as Code (IaC) sécurisée.} La configuration
    des infrastructures est déclarée sous forme de code versionné, soumis aux
    mêmes revues et analyses de sécurité que le code applicatif.

\end{enumerate}

% ------------------------------------------------------------
%  1.3  LES TYPES DE TESTS DE SÉCURITÉ
% ------------------------------------------------------------
\section{Les types de tests de sécurité (SAST, DAST, IAST, RASP)}
\label{sec:tests}

La sécurité applicative mobilise plusieurs familles de techniques de test,
chacune offrant un angle d'observation distinct sur la posture de sécurité
d'un système. Nous présentons ici les quatre approches majeures.

% ── 1.3.1  SAST ─────────────────────────────────────────────
\subsection{SAST  Static Application Security Testing}

Le SAST (test de sécurité statique) analyse le code source, le bytecode ou
le code binaire d'une application \textbf{sans l'exécuter}, afin
d'identifier des vulnérabilités telles que les injections SQL, les
dépassements de tampon, ou les mauvaises pratiques cryptographiques.


\textbf{Principes de fonctionnement.} Le SAST opère par analyse syntaxique et
sémantique du code. Il construit généralement un \emph{Abstract Syntax Tree}
(AST) et effectue une analyse du flux de données (\emph{taint analysis}) pour
tracer le cheminement des données non fiables depuis leurs sources jusqu'aux
points d'utilisation potentiellement dangereux (\emph{sinks}).

\textbf{Avantages.}
\begin{itemize}[leftmargin=*]
  \item Détection précoce des vulnérabilités, dès la phase de codage.
  \item Couverture exhaustive du code, y compris les chemins d'exécution
    rarement empruntés.
  \item Intégration aisée dans les IDE et les pipelines CI.
\end{itemize}

\textbf{Limites.}
\begin{itemize}[leftmargin=*]
  \item Taux élevé de faux positifs, nécessitant une validation humaine.
  \item Incapacité à détecter les vulnérabilités de configuration ou les
    failles logiques dépendant du contexte d'exécution.
  \item Difficulté à analyser les dépendances tierces compilées.
\end{itemize}

\textbf{Outils représentatifs.} SonarQube, Semgrep, Checkmarx, Fortify,
Bandit (Python), SpotBugs (Java).

% ── 1.3.2  DAST ─────────────────────────────────────────────
\subsection{DAST  Dynamic Application Security Testing}


Le DAST (test de sécurité dynamique) évalue la sécurité d'une application
\textbf{en cours d'exécution}, en simulant des attaques réelles depuis
l'extérieur, sans accès au code source. Il adopte la perspective d'un
attaquant (\emph{black-box testing}).


\textbf{Principes de fonctionnement.} Le DAST interagit avec l'application
via ses interfaces exposées (API REST, interfaces web, gRPC, etc.), en
injectant des charges malveillantes (\emph{payloads}) et en analysant les
réponses pour détecter des comportements anormaux révélateurs de
vulnérabilités~: injections, authentification défaillante, exposition de
données sensibles, etc.

\textbf{Avantages.}
\begin{itemize}[leftmargin=*]
  \item Indépendant du langage et du framework applicatif.
  \item Détecte des vulnérabilités réelles dans l'environnement d'exécution,
    incluant celles issues de la configuration et des interactions entre
    composants.
  \item Faible taux de faux positifs par rapport au SAST.
\end{itemize}

\textbf{Limites.}
\begin{itemize}[leftmargin=*]
  \item Couverture partielle~: ne peut tester que les fonctionnalités
    accessibles et les surfaces d'attaque exposées.
  \item Détection tardive dans le cycle de développement.
  \item Complexité accrue dans les architectures microservices, où la
    multitude des points d'entrée démultiplie la surface à couvrir.
\end{itemize}

\textbf{Outils représentatifs.} OWASP ZAP, Burp Suite, Nikto, Nuclei,
Nessus.

% ── 1.3.3  IAST ─────────────────────────────────────────────
\subsection{IAST  Interactive Application Security Testing}


  L'IAST (test de sécurité interactif) \textbf{combine les approches statique
  et dynamique} en instrumentant l'application à l'aide d'agents déployés
  dans son environnement d'exécution. Il analyse le comportement interne du
  code pendant les tests fonctionnels.


\textbf{Principes de fonctionnement.} Des agents IAST sont intégrés
directement dans le runtime de l'application (JVM, CLR, interpréteur Python,
etc.). Ils interceptent les flux d'exécution, tracent les données et
détectent les vulnérabilités en observant les interactions réelles entre les
composants, sans rejeter les requêtes légitimes.

\textbf{Avantages.}
\begin{itemize}[leftmargin=*]
  \item Très faible taux de faux positifs grâce à la corrélation contexte
    d'exécution\,/\,code source.
  \item Fonctionnement en temps réel, sans phase de scan dédiée.
  \item Adapté aux environnements de test d'intégration continue.
\end{itemize}

\textbf{Limites.}
\begin{itemize}[leftmargin=*]
  \item Surcharge de performance liée à l'instrumentation.
  \item Support limité à certains langages et frameworks.
  \item Couverture dépendante de l'exhaustivité des tests fonctionnels
    exécutés.
\end{itemize}

\textbf{Outils représentatifs.} Contrast Security, Seeker (Synopsys), Hdiv
Security.

% ── 1.3.4  RASP ─────────────────────────────────────────────
\subsection{RASP  Runtime Application Self-Protection}


  Le RASP (auto-protection applicative à l'exécution) est une technologie de
  sécurité \textbf{intégrée directement dans une application} ou son
  environnement d'exécution, capable de détecter et de bloquer des attaques
  en temps réel, sans intervention externe.


\textbf{Principes de fonctionnement.} Contrairement au DAST qui teste de
l'extérieur, le RASP s'exécute à l'intérieur de l'application. Il intercepte
les appels système, les requêtes de base de données et autres opérations
sensibles, et peut les bloquer instantanément si un comportement malveillant
est détecté, tout en émettant des alertes vers le SIEM.

\textbf{Avantages.}
\begin{itemize}[leftmargin=*]
  \item Protection en production, y compris contre des attaques zero-day.
  \item Réduction de la dépendance aux mises à jour du WAF.
  \item Contexte d'exécution précis, minimisant les faux positifs.
\end{itemize}

\textbf{Limites.}
\begin{itemize}[leftmargin=*]
  \item Impact sur les performances applicatives.
  \item Risque de faux positifs entraînant des interruptions de service.
  \item Complexité de déploiement et de maintenance.
\end{itemize}

\textbf{Outils représentatifs.} Sqreen (Datadog), Imperva RASP, OpenRASP.

% ── 1.3.5  Tableau comparatif ───────────────────────────────
\subsection{Tableau comparatif des quatre approches}

\begin{table}[h!]
\centering
\caption{Comparaison des approches de tests de sécurité}
\label{tab:comparaison-tests}
\renewcommand{\arraystretch}{1.3}
\begin{tabularx}{\textwidth}{>{\bfseries}lXXXX}
\toprule
\textbf{Critère} & \textbf{SAST} & \textbf{DAST} & \textbf{IAST} & \textbf{RASP} \\
\midrule
Phase           & Codage / build    & Test / staging  & Test / intég.   & Production \\
Accès code      & Requis            & Non requis      & Optionnel       & Non requis \\
Perspective     & White-box         & Black-box       & Gray-box        & Interne \\
Faux positifs   & Élevés            & Faibles         & Très faibles    & Faibles \\
Couverture      & Code complet      & Surface exposée & Code exercé     & Runtime \\
Impact perf.    & Aucun             & Faible          & Moyen           & Moyen--élevé \\
Microservices   & Partiel           & Complexe        & Adapté          & Adapté \\
\bottomrule
\end{tabularx}
\end{table}

% ------------------------------------------------------------
%  1.4  FOCUS DAST : ENJEUX MICROSERVICES
% ------------------------------------------------------------
\section{Focus DAST : enjeux dans les architectures microservices}
\label{sec:dast-micro}

\subsection{Spécificités des architectures microservices}

Une architecture microservices est un style architectural dans lequel une
application est décomposée en un ensemble de services petits, autonomes et
faiblement couplés, chacun responsable d'une capacité métier précise. Ces
services communiquent entre eux via des API bien définies, généralement sur
des protocoles HTTP/REST ou gRPC, et peuvent être déployés indépendamment.

Cette architecture présente plusieurs caractéristiques qui complexifient
significativement l'application du DAST~:

\paragraph{Explosion de la surface d'attaque.}
Dans une application monolithique, la surface d'attaque exposée est
concentrée sur un nombre limité de points d'entrée. Dans une architecture
microservices, chaque service expose potentiellement sa propre API,
multipliant d'autant les surfaces à tester. Un système composé de
50~microservices peut exposer plusieurs centaines d'endpoints distincts, avec
des contrats d'API hétérogènes.

\paragraph{Dynamicité de l'infrastructure.}
Les microservices s'exécutent dans des environnements hautement dynamiques
(Kubernetes, par exemple), où les instances de services apparaissent et
disparaissent selon la charge. Les adresses IP et les ports sont attribués
dynamiquement, rendant obsolète toute cible DAST définie statiquement. Un
scanner DAST doit s'adapter en temps réel à cette topologie mouvante.

\paragraph{Communications interservices non exposées.}
Une partie significative des communications dans une architecture
microservices est interne~: les échanges entre services du backend ne sont
pas directement accessibles depuis l'extérieur du cluster. Les outils DAST
conventionnels, conçus pour tester des interfaces web publiques, ne peuvent
pas atteindre ces flux de communication internes sans mécanismes spécifiques.

\paragraph{Hétérogénéité des protocoles et des données.}
Les microservices peuvent utiliser des protocoles variés~: REST/HTTP, gRPC,
GraphQL, WebSocket, messaging asynchrone (Kafka, RabbitMQ). Les outils DAST
doivent être capables de générer des payloads adaptés à chaque protocole et
d'interpréter correctement les réponses correspondantes.

\subsection{Stratégies DAST adaptées aux microservices}

Face à ces défis, plusieurs stratégies permettent d'adapter le DAST aux
architectures microservices~:

\textbf{Intégration dans la pipeline CI/CD.} Plutôt que d'effectuer des
scans ponctuels, le DAST est intégré comme une étape automatique de la
pipeline, déclenchée à chaque déploiement en environnement de staging. Des
outils comme OWASP ZAP disposent de modes \og{}headless\fg{} permettant cette
intégration.

\textbf{DAST orienté API.} Des outils spécialisés dans le test de sécurité
des API (comme 42Crunch, Traceable~AI, ou les extensions REST/gRPC de ZAP)
permettent d'importer les spécifications OpenAPI/Swagger et de générer
automatiquement des cas de test de sécurité couvrant l'ensemble des endpoints
documentés.

\textbf{Service Mesh comme point d'observation.} L'adoption d'un Service Mesh
ouvre une perspective nouvelle pour le DAST dans les microservices. En
opérant au niveau du sidecar proxy, le Service Mesh fournit une visibilité
totale sur les communications interservices, y compris les flux internes au
cluster, sans modification du code applicatif. Cette caractéristique permet
d'envisager un DAST interne, capable de tester la sécurité des échanges entre
services~--- une dimension absente des approches DAST traditionnelles.


\textbf{Synergie clé~:} Le Service Mesh ne se substitue pas au DAST, mais
il en élargit considérablement le périmètre d'action dans les architectures
microservices. En fournissant un point d'interception centralisé et une
observabilité totale du trafic interservices, il crée les conditions d'un
DAST interne couvrant les vulnérabilités invisibles depuis l'extérieur du
cluster.


\subsection{Limites actuelles et pistes d'amélioration}

Malgré les évolutions récentes, plusieurs limites persistent dans
l'application du DAST aux architectures microservices~:

\begin{itemize}[leftmargin=*]
  \item \textbf{Couverture incomplète des flux asynchrones~:} les
    communications via des bus de messages (Kafka, RabbitMQ) restent
    difficiles à intercepter et à tester dynamiquement.
  \item \textbf{Manque de standardisation~:} l'absence d'un format universel
    de description des API interservices limite la capacité des outils DAST à
    générer automatiquement des cas de test pertinents.
  \item \textbf{Complexité de l'environnement de test~:} reproduire
    fidèlement un environnement microservices complet à des fins de test reste
    coûteux en ressources et en configuration.
\end{itemize}

Ces limites motivent l'exploration d'une approche hybride intégrant le
Service Mesh comme infrastructure de sécurité complémentaire, objet des
chapitres suivants.

% ------------------------------------------------------------
%  1.5  CONCLUSION
% ------------------------------------------------------------
\section{Conclusion}
\label{sec:conclusion-ch1}

Ce chapitre a posé les fondements conceptuels et méthodologiques sur lesquels
repose ce projet. Nous avons d'abord formulé les objectifs du projet en trois
axes~: cadre méthodologique DevSecOps, architecture hybride proposée et
validation expérimentale. Nous avons ensuite introduit le paradigme DevSecOps,
en retraçant son émergence depuis le mouvement DevOps et en détaillant ses
principes fondateurs. Nous avons dressé un panorama structuré des quatre
grandes familles de tests de sécurité SAST, DAST, IAST et RASP en
mettant en évidence leurs complémentarités et leurs limites respectives.

La section finale a approfondi les enjeux spécifiques du DAST dans les
architectures microservices, en identifiant les défis majeurs que soulève
cette combinaison~: explosion de la surface d'attaque, dynamicité de
l'infrastructure, inaccessibilité des flux internes et hétérogénéité
protocolaire.



